\section{抽象表示：医学推理与多语言回路}

方法建立后，第一个发现是模型内部存在跨表面形式复用的抽象表示。本章用医学与多语言两个案例说明：同一个概念可以被不同词语激活，而输入语言和输出语言的边界信息又可以与中间语义部分分离。

\subsection{医学案例：症状、诊断与判据在一遍前向传播中组合}

病例给出孕晚期、右上腹痛、高血压和肝酶升高，要求只问一个额外症状。模型回答 visual disturbances。重要问题不是医学结论本身，而是图中是否出现 pregnancy、preeclampsia、diagnostic criteria 等可组合表示，以及压制某个 feature 后答案是否发生预期改变。

\lecturefigure{slide-25-medical-reasoning-prompt.jpg}{单次前向传播中的医学推理 prompt}{00:37:30}

\lecturefigure{slide-26-medical-concept-representations.jpg}{模型把症状、诊断与额外诊断判据组织为抽象 feature}{00:38:46}

图 26 先由多个症状 feature 汇聚到疾病相关表示，再由诊断表示推动“应询问 visual disturbances”。这不是一条自然语言 chain of thought，却体现了结构化关系：症状不是逐字匹配，而是作为诊断证据在内部汇合。

\lecturefigure{slide-27-medical-feature-interventions.jpg}{替换或抑制医学 feature 后，推荐追问随之改变}{00:39:51}

干预提高了因果可信度：如果把 biliary-system 相关 feature 推高，输出向另一类症状移动，就说明这些方向不只是事后命名。不过临床知识可能不完整、prompt 也很简化；图不能证明模型适合真实诊断，更不能替代医学指南与医生判断。

\begin{warningbox}{医学案例的范围}
这是 mechanistic case study，不是医疗建议或临床可靠性证明。它说明模型能在一个 prompt 上组合医学表示；不说明知识完整、概率校准、群体公平或真实场景安全。
\end{warningbox}

\teachervoice{00:31:52--00:42:44，老师把医学例子作为“一次 forward pass 如何完成组合推理”的窗口，并在问答中承认真实医学侧的可靠性需要独立验证。}

\subsection{多语言案例：边界是语言，核心是共享语义}

“small 的反义词”可用英语、法语或中文提问。若模型完全在英语中思考，整个内部图应先翻译成英语再计算；若每种语言完全独立，三套图几乎不共享。实际结果更像分层结构：输入和输出附近有语言特定 feature，中间则出现 multilingual small、antonym、large 等共享表示。

\lecturefigure{slide-28-unique-language-question.jpg}{问题：模型是否在某一种唯一语言中“思考”}{00:42:47}

\lecturefigure{slide-29-shared-multilingual-representations.jpg}{三种语言共享中间语义 feature，同时保留语言特定边界}{00:45:01}

读图时先看两端：quote、English/Chinese/French 与具体 token 负责语言形式；再看中间：small、opposite、say large 等 feature 跨语言重合。结论不是“模型有纯粹思想语言”，而是语义抽象与表面实现部分解耦。

\lecturefigure{slide-30-antonym-to-synonym-edit.jpg}{把 antonym feature 换成 synonym feature，三种语言的输出都按语义改变}{00:45:06}

这项 intervention 很关键：同一内部编辑能跨语言把“反义词”任务改成“同义词”任务，说明共享 feature 参与运算而非只与输出相关。输入的 operand、操作类型和输出语言可以分别编辑，展示了模块化但非完全离散的内部表示。

若两种语言的活跃 feature 集分别为 $F_a,F_b$，共享比例可用交并比表示：
\[
\operatorname{IoU}(F_a,F_b)=\frac{|F_a\cap F_b|}{|F_a\cup F_b|}.
\]
其中交集计数共同 feature，并集计数至少在一种语言中出现的 feature。IoU 上升表示共享更多，但受 feature matching、阈值和替换模型质量影响。

\lecturefigure{slide-31-representation-sharing-scales.jpg}{更大模型在语言之间共享更高比例的表示}{00:47:26}

图 31 的趋势支持“规模增大促进抽象复用”，但不证明所有概念都会收敛到同一语义空间。语言距离、tokenization、数据量和文化特定表达仍会造成差异；曲线是该实验设置下的 feature overlap，不是普适语言学定律。

\begin{knowledgebox}{从“思考语言”改写为三个可测问题}
不要问模型到底“用哪种语言想”。可以分别问：输入阶段哪些 feature 依赖语言形式；中间操作是否共享语义方向；输出阶段如何选择目标语言。三个问题能被图和干预检验，单一拟人化问题则容易混淆。
\end{knowledgebox}

\subsection{本章小结}

医学与多语言案例共同表明，模型会把表面 token 组合成更抽象的概念和操作。抽象表示可跨症状、语言与表达复用，但仍嵌在具体 prompt、模型和 feature 提取方法中；“可复用”不能升级为“普遍、唯一、完整”。

